\documentclass[11pt]{article}

\usepackage[margin=1in]{geometry}
\usepackage{amsmath, amssymb, amsthm}
\usepackage{enumitem}

\pagestyle{empty}

\begin{document}

\begin{center}
  {\Large\bfseries Weekly Assignment: Logic\\[2pt]
   (Epp~2.1 \& 2.2)}
\end{center}

\bigskip

\noindent
Complete each of the following problems.  Show all work and justify your answers.

\bigskip

%% ─────────────────────────────────────────────
\noindent\textbf{Problem 1} (Epp 2.1 \#4). \;
Represent the common form of each argument using letters to
stand for component sentences, and fill in the blanks so that the
argument in part~(b) has the same logical form as the argument
in part~(a).

\begin{enumerate}[label=\textbf{\alph*.},itemsep=6pt]
  \item If the program syntax is faulty, then the computer will
        generate an error message.\\
        If the computer generates an error message, then the
        program will not run.\\
        Therefore, if the program syntax is faulty, then the
        program will not run.

  \item If this simple graph \rule{2cm}{0.4pt}\,, then it is complete.\\
        If this graph \rule{2cm}{0.4pt}\,, then any two of its
        vertices can be joined by a path.\\
        Therefore, if this simple graph has 4 vertices and 6 edges,
        then \rule{2cm}{0.4pt}\,.
\end{enumerate}

\bigskip

%% ─────────────────────────────────────────────
\noindent\textbf{Problem 2} (Epp 2.1 \#8\,b,c,e). \;
Write the following statements in symbolic form using the
symbols $\lnot$, $\lor$, and~$\land$.  Let $h ={}$``John is
healthy,'' $w ={}$``John is wealthy,'' and $s ={}$``John is
wise.''

\begin{enumerate}[label=\textbf{\alph*.},start=2,itemsep=4pt]
  \item John is not wealthy but he is healthy and wise.
  \setcounter{enumi}{2}
  \item John is neither healthy, wealthy, nor wise.
  \setcounter{enumi}{4}
  \item John is wealthy, but he is not both healthy and wise.
\end{enumerate}

\bigskip

%% ─────────────────────────────────────────────
\noindent\textbf{Problem 3} (Epp 2.1 \#9\,b,c). \;
Write the following statements in symbolic form using the
symbols $\lnot$, $\lor$, and~$\land$.  Let $p ={}$``$x > 5$,''
$q ={}$``$x = 5$,'' and $r ={}$``$10 > x$.''

\begin{enumerate}[label=\textbf{\alph*.},start=2,itemsep=4pt]
  \item $10 > x > 5$
  \setcounter{enumi}{2}
  \item $10 > x \ge 5$
\end{enumerate}

\bigskip

%% ─────────────────────────────────────────────
\noindent\textbf{Problem 4} (Epp 2.1 \#22). \;
Determine whether the following statement forms are logically
equivalent.  Construct a truth table and include a sentence
justifying your answer.
\[
  p \land (q \lor r)
  \qquad\text{and}\qquad
  (p \land q) \lor (p \land r)
\]

\bigskip

%% ─────────────────────────────────────────────
\noindent\textbf{Problem 5} (Epp 2.1 \#24). \;
Determine whether the following statement forms are logically
equivalent.  Construct a truth table and include a sentence
justifying your answer.
\[
  (p \lor q) \lor (p \land r)
  \qquad\text{and}\qquad
  (p \lor q) \land r
\]

\bigskip

%% ─────────────────────────────────────────────
\noindent\textbf{Problem 6} (Epp 2.1 \#42). \;
Use a truth table to determine whether the following statement
form is a tautology, a contradiction, or neither.
\[
  \bigl((\lnot p \land q) \land (q \land r)\bigr) \land \lnot q
\]

\bigskip

%% ─────────────────────────────────────────────
\noindent\textbf{Problem 7} (Epp 2.2 \#2). \;
Rewrite the following statement in if-then form:

\medskip
\begin{quote}
  \emph{I am on time for work if I catch the 8:05 bus.}
\end{quote}

\bigskip

%% ─────────────────────────────────────────────
\noindent\textbf{Problem 8} (Epp 2.2 \#18). \;
Write each of the following three statements in symbolic form
and determine which pairs are logically equivalent.  Include
truth tables and a few words of explanation.

\begin{enumerate}[label=(\roman*),itemsep=4pt]
  \item If it walks like a duck and it talks like a duck,
        then it is a duck.
  \item Either it does not walk like a duck or it does not talk
        like a duck, or it is a duck.
  \item If it does not walk like a duck and it does not talk like
        a duck, then it is not a duck.
\end{enumerate}

\end{document}
